\documentclass[12pt]{article}
\usepackage[T1]{fontenc}
\usepackage{lmodern}
\usepackage{amsmath,amsfonts,amssymb,amsthm}
\usepackage[shortlabels]{enumitem}
\usepackage{mathtools}
\usepackage{fullpage}
\usepackage{microtype}
\usepackage{needspace}
\usepackage{booktabs,tabularx}
\usepackage[hidelinks]{hyperref}
\hypersetup{pdftitle={Algebraic rootfinding methods: a conceptual survey},
  pdfauthor={Etna Lindy, Martin Lotz, Vanni Noferini}}
\setlength{\emergencystretch}{2em}

\def\C{\mathbb{C}}
\def\R{\mathbb{R}}
\def\B{\mathcal{B}}
\def\Scal{\mathcal{S}}
\DeclareMathOperator{\Res}{Res}
\DeclareMathOperator{\adj}{adj}
\DeclareMathOperator{\tr}{trace}
\DeclareMathOperator{\lc}{lc}
\DeclareMathOperator{\rank}{rank}
\DeclareMathOperator{\diag}{diag}
\DeclareMathOperator{\spanop}{span}
\DeclareMathOperator{\NF}{NF}

\newtheorem{theorem}{Theorem}[section]
\newtheorem{proposition}[theorem]{Proposition}
\newtheorem{lemma}[theorem]{Lemma}
\newtheorem{corollary}[theorem]{Corollary}
\theoremstyle{definition}
\newtheorem{definition}[theorem]{Definition}
\newtheorem{example}[theorem]{Example}
\theoremstyle{remark}
\newtheorem{remark}[theorem]{Remark}
\newenvironment{reduction}
  {\par\medskip\noindent\textbf{Reduction and recovery.}\ }
  {\par\medskip}

\title{Algebraic rootfinding methods: a conceptual survey}
\author{Etna Lindy, Martin Lotz, Vanni Noferini}
\date{}

\begin{document}
\maketitle

\begin{abstract}
Algebraic rootfinders represent the finite zero set of a polynomial system by a matrix,
a matrix polynomial, or a univariate polynomial. We explain seven constructions:
Sylvester and Cayley--B\'ezout resultants, Macaulay matrices, multiparameter eigenvalue
problems, Gr\"obner elimination, rational univariate representations, and normal-form
methods. A common quotient-algebra framework explains why these representations work,
while a single four-root example makes their differences explicit. Short proofs cover
the main algebraic identities; more substantial existence and regularity results are
stated with their hypotheses and references. We distinguish exact root correspondence
from numerical stability, including the effects of projection, basis choice, and the
norm in which data are perturbed. The survey provides a methodological foundation for
studying conditioning without presupposing an average-case verdict on the methods.
\end{abstract}

\section{The problem and a guide to the methods}\label{sec:intro}

Let $P=(p_1,\ldots,p_d)$ be a polynomial map $\C^d\to\C^d$. We assume throughout that
$I=\langle p_1,\ldots,p_d\rangle$ is radical and zero-dimensional, with distinct roots
$z_1,\ldots,z_N$, and that $J(z_j)=DP(z_j)$ is invertible. These assumptions exclude
multiple roots and positive-dimensional solution sets. They make the central ideas
transparent; methods for multiplicities require additional local-algebra machinery
\cite{moller1995multivariate,stetter2004numerical}.

The task is to compute all $N$ roots. Newton iteration is local, whereas homotopy
continuation can find all roots by tracking a suitable collection of paths.
The methods considered here organize the computation differently:
\[
\text{coefficients}\ \longrightarrow\ \text{algebraic representation}
\ \longrightarrow\ \text{spectral or scalar roots}
\ \longrightarrow\ \text{solutions of }P.
\]
We will ask three questions of each construction: what is built, why does it encode
the roots, and how are all coordinates recovered? The answers also identify three
places where numerical errors can enter: construction, solution of the intermediate
problem, and recovery. Exact equivalence of representations does not imply equal
conditioning of these operations.

\Needspace{7\baselineskip}
\subsection{Three algorithmic families}\label{sec:taxonomy}

\begin{description}[style=nextline]
  \item[Resultant and eigenvalue constructions.]
  Build a matrix that becomes singular at a root or at a projected root.
  Sylvester and Cayley--B\'ezout methods hide one variable; Macaulay methods adjoin
  a linear form; multiparameter methods replace the equations by affine matrix pencils.
  \item[Scalar elimination.]
  Build a univariate polynomial whose roots label the solutions, together with
  formulas recovering their coordinates. Gr\"obner elimination and rational univariate
  representations take this route.
  \item[Normal forms.]
  Construct multiplication operators on the finite-dimensional quotient algebra
  $A=\C[x_1,\ldots,x_d]/I$, then recover the roots from their common eigenvectors.
\end{description}

These families overlap. A scalar eliminant can be a characteristic polynomial of a
multiplication matrix, and a Macaulay matrix can be used to construct a normal form.
The classification describes how the representation is computed, not mutually exclusive
algebraic objects. Table~\ref{tab:methods} at the end collects their sizes and hypotheses.

\begin{example}[A running system]\label{ex:running}
We will follow the same system through the constructions:
\begin{equation}\label{eq:running}
  p(x,y)=x^2-1=0,\qquad q(x,y)=y^2-x=0.
\end{equation}
Its roots are
\[
  (1,1),\quad(1,-1),\quad(-1,i),\quad(-1,-i).
\]
Indeed $x=y^2$ and $y^4=1$. The Jacobian determinant is $4xy$, which is nonzero
at all four roots. The coordinate $y$ distinguishes the roots, while $x$ does not:
there are two roots above each value of $x$. This distinction will explain why a
simple polynomial root need not give a simple projected eigenvalue.
\end{example}

\subsection{Degree conventions and roots at infinity}\label{sec:degrees}

For total degrees $e_1,\ldots,e_d$, a generic system has $\prod_i e_i$ projective
roots by B\'ezout's theorem, and all are affine if there is no common root on the
hyperplane at infinity. If instead every equation has degree at most $n$ in
\emph{each} variable, its support lies in the box $[0,n]^d$. The generic count is
$d!n^d$, naturally computed in $(\mathbb P^1)^d$; see the multihomogeneous and sparse
resultant accounts in \cite{cox2005using,gelfand1994discriminants}. The running example
has total degrees $(2,2)$ and four roots; it is not a generic box-degree-two system,
which would have eight.

The compactification matters. Ordinary homogenization of a full box-degree-$n$ equation
to total degree $dn$ has top part $c(x_1\cdots x_d)^n$. For $d\ge2$, the resulting
equations have common zeros at ordinary projective infinity even when the system is
generic in its box support. Generic absence of boundary zeros in $(\mathbb P^1)^d$
does not remove these ordinary projective zeros. We therefore state the no-infinity
assumption separately for the dense Macaulay construction.

\section{The common algebra: functions on a finite root set}\label{sec:algebra}

Write $R=\C[x_1,\ldots,x_d]$ and $A=R/I$. The class $[f]\in A$ records the values
of $f$ at the roots; adding a polynomial in $I$ changes none of them. The following
elementary form of the Chinese remainder theorem is the organizing fact.

\begin{proposition}[Evaluation and interpolation]\label{prop:evaluation}
Evaluation is an algebra isomorphism
\[
  \operatorname{ev}:A\longrightarrow\C^N,\qquad
  [f]\longmapsto(f(z_1),\ldots,f(z_N)),
\]
where multiplication on $\C^N$ is componentwise. In particular $\dim A=N$.
There are unique classes $\ell_j\in A$ with $\ell_j(z_k)=\delta_{jk}$;
they satisfy $\ell_j\ell_k=\delta_{jk}\ell_j$ and $\sum_j\ell_j=1$.
\end{proposition}

\begin{proof}
Because $I$ is radical, a polynomial vanishes at every root exactly when it belongs
to $I$ (Hilbert's Nullstellensatz), so evaluation is injective.
Choose a linear form $t$ taking distinct values $t_j=t(z_j)$. Such choices exist:
the prohibited coefficient vectors lie in the finite union of hyperplanes
$\{u:u^T(z_j-z_k)=0\}$, $j\ne k$. The polynomials
\[
  \ell_j(x)=\prod_{k\ne j}\frac{t(x)-t_k}{t_j-t_k}
\]
evaluate to the coordinate vectors of $\C^N$, proving surjectivity. The remaining
identities follow by evaluation and injectivity.
\end{proof}

A form taking distinct values on the root set is called \emph{separating}.
The proof is conceptual: algorithms compute the quotient from the coefficients;
they do not first compute the roots to construct the $\ell_j$.

\subsection{Multiplication becomes diagonal}\label{sec:multiplication}

For $g\in R$, multiplication by $[g]$ is a linear operator on $A$. In a basis
$\B=(b_1,\ldots,b_N)$, write its matrix as $M_g$, with column $k$ equal to
$[gb_k]_\B$. Define the evaluation matrix
\[
  V=(b_k(z_j))_{k,j},\qquad v_j=\B(z_j).
\]
Columns index roots; coefficient vectors are columns. Throughout, ${}^T$ denotes
transpose and ${}^*$ conjugate transpose.

\begin{proposition}[The multiplication eigenvalue theorem]\label{prop:multiplication}
The matrix $V$ is invertible and
\begin{equation}\label{eq:stickelberger}
  M_g=V^{-T}\diag(g(z_1),\ldots,g(z_N))V^T.
\end{equation}
Thus all multiplication matrices commute and are simultaneously diagonalizable.
Their characteristic and minimal polynomials are, respectively,
\[
  \det(TI-M_g)=\prod_{j=1}^N(T-g(z_j)),\qquad
  \prod_{\eta\in\{g(z_1),\ldots,g(z_N)\}}(T-\eta).
\]
\end{proposition}

\begin{proof}
Evaluation sends the coefficient vector $c$ to $V^Tc$ and sends multiplication by $g$
to multiplication by $\diag(g(z_j))$. Proposition~\ref{prop:evaluation} gives the
invertibility of $V$ and the asserted similarity. The polynomial identities follow
for a diagonal matrix.
\end{proof}

In this convention the right eigenvector at a separated value $g(z_j)$ is
$[\ell_j]_\B=V^{-T}e_j$, and $v_j^T$ is a left eigenvector. Equivalently,
the evaluation vectors are right eigenvectors of $M_g^T$. Repeated values of $g$
create higher-dimensional eigenspaces even when every root of $P$ is simple.

\subsection{The univariate model}\label{sec:companion}

For a monic $p(T)=T^m+a_{m-1}T^{m-1}+\cdots+a_0$, reduction modulo $p$ in the
basis $(1,T,\ldots,T^{m-1})$ gives the companion matrix
\begin{equation}\label{eq:companion}
  C_p=\begin{pmatrix}
  0&&&-a_0\\
  1&0&&-a_1\\
  &\ddots&\ddots&\vdots\\
  &&1&-a_{m-1}
  \end{pmatrix},\qquad \det(TI-C_p)=p(T).
\end{equation}
For squarefree $p$ the determinant identity follows from
Proposition~\ref{prop:multiplication}; for every monic $p$ it follows because
$e_1,C_pe_1,\ldots,C_p^{m-1}e_1$ is a basis and $p(C_p)=0$, so the minimal
and characteristic polynomials both equal $p$. At a root $\lambda$,
$(1,\lambda,\ldots,\lambda^{m-1})$ is a \emph{left} eigenvector.

This is the basic representation change: a scalar polynomial equation becomes a
matrix eigenproblem. The multivariate methods make different choices about which
part of this pattern to preserve.

\section{Family I: resultant and eigenvalue constructions}\label{sec:familyI}

Here singular matrices detect roots. A \emph{regular} square matrix polynomial
$R(\lambda)$ has determinant not identically zero. Its finite eigenvalues are the
zeros of that determinant; a singular matrix polynomial requires a different analysis.
Even in the regular case, degree drops can produce eigenvalues at infinity. The
following constructions specify when finite eigenvalues represent affine roots.
\subsection{The Sylvester resultant (\texorpdfstring{$d=2$}{d=2})}\label{sec:sylvester}

\emph{Idea.} Freeze one coordinate, say $y$, and ask whether the resulting
polynomials in $x$ have a common root. The Sylvester matrix turns this question
into singularity of a matrix whose entries are polynomials in $y$.

\emph{Construction and proof.} First take nonzero univariate polynomials
$f(x)=a_mx^m+\cdots+a_0$ and $g(x)=b_kx^k+\cdots+b_0$ of
\emph{actual} degrees $m,k\ge1$. Write the coefficient vectors of
\[
  x^{k-1}f,\ldots,f,\quad x^{m-1}g,\ldots,g
\]
as rows, with columns ordered $x^{m+k-1},\ldots,1$. This defines the
$(m+k)\times(m+k)$ Sylvester matrix $S(f,g)$. Its transpose represents the map
\[
  \C[x]_{<k}\oplus\C[x]_{<m}\longrightarrow\C[x]_{<m+k},
  \qquad (u,v)\longmapsto uf+vg.
\]
If $f$ and $g$ are coprime, $uf=-vg$ implies that $f$ divides $v$; the degree
bound forces $v=0$, then $u=0$. Conversely, if $h=\gcd(f,g)$ has positive
degree, $(u,v)=(g/h,-f/h)$ is a nonzero vector in this kernel. Thus
\[
  S(f,g)\text{ is singular}\quad\Longleftrightarrow\quad
  f,g\text{ have a common root in }\C.
\]
With this ordering its determinant is the resultant, with product formula
\[
  \det S(f,g)=\Res_x(f,g)
  =a_m^{\,k}\prod_{f(\alpha)=0}g(\alpha),
\]
where roots are counted with multiplicity; see \cite[Ch.~3]{cox2005using}.
The same argument shows $\dim\ker S(f,g)=\deg\gcd(f,g)$: every relation is
$(u,v)=w(g/h,-f/h)$ with $\deg w<\deg h$.

For $p,q\in\C[x,y]$, choose their degrees $m,k$ in $x$ over $\C[y]$ and
form $S(y)$ by the same shifts. Its determinant is a polynomial in $y$.
At a parameter value where both degrees are preserved, the preceding proof
applies verbatim. If $(x_*,y_*)$ is an affine common root, direct evaluation
of the rows gives
\[
  S(y_*)\begin{pmatrix}x_*^{m+k-1}&\cdots&x_*&1\end{pmatrix}^{T}=0.
\]
When the specialized gcd is linear, this vector spans the right kernel, so
its last two entries recover $x_*$ by a ratio. Several roots in the same
fiber require additional work, such as computing the fiber gcd.

\begin{example}[The running system]\label{ex:sylvester-running}
  For $p=x^2-1$ and $q=y^2-x$ in Example~\ref{ex:running}, the degrees in
  $x$ are $2$ and $1$. The rows $p,xq,q$ give
  \[
    S(y)=\begin{pmatrix}1&0&-1\\-1&y^2&0\\0&-1&y^2\end{pmatrix},
    \qquad \det S(y)=y^4-1.
  \]
  At each $\eta\in\{1,-1,\mathrm{i},-\mathrm{i}\}$, a right null vector is
  $(\eta^4,\eta^2,1)^T=(1,\eta^2,1)^T$. Thus $x=\eta^2$, recovering
  $(1,\pm1)$ and $(-1,\pm\mathrm{i})$. All four $y$-values are distinct;
  hiding $x$ instead would combine the roots into two fibers.
\end{example}

\begin{remark}[A resultant remembers its degree convention]\label{rem:sylvester-infinity}
  A matrix constructed with fixed degrees remains that size when leading
  coefficients vanish after specialization. It then tests common zeros of
  the \emph{homogenized} polynomials in $\mathbb P^1$, including infinity.
  For example, $f=ax+1$, $g=ax+2$ give
  \[
    S(a)=\begin{pmatrix}a&1\\a&2\end{pmatrix},\qquad\det S(a)=a.
  \]
  At $a=0$ the affine polynomials $1$ and $2$ have no root, but their
  degree-one homogenizations $aX+Z$, $aX+2Z$ share $[X:Z]=[1:0]$.
  Consequently a vanishing fixed-degree resultant need not identify an
  affine solution. Degree drops and boundary roots must be checked when
  extending a projected candidate to a full solution.
\end{remark}

\begin{reduction}
  Solve the polynomial eigenvalue problem $S(y)v=0$. For generic polynomials
  of degree $n$ in each variable, $S$ has size $2n$, degree $n$, and
  $2n^2$ finite eigenvalues, one per root. A companion linearization has
  size $2n^2$ and can be solved by QZ. Sizes and grades must be adapted to
  other supports: the running example has size $3$ and degree $2$, so its
  standard size-$6$ linearization also has two eigenvalues at infinity.
\end{reduction}

\emph{Numerical aspects and references.} The entries of $S$ depend linearly
on the input coefficients. Nevertheless, arbitrary perturbations of $S$'s
coefficient matrices need not respect the repeated-coefficient structure
of a Sylvester matrix. QZ controls backward error for its pencil; translating
this into error in the original polynomials also requires control of the
linearization and construction. See \cite{tisseur2000backward,
noferini2016numerical} for these numerical issues. The classical sources are
\cite{sylvester1840method,sylvester1853theory}; accessible accounts include
\cite[Ch.~3]{cox2005using} and \cite{sturmfels2002solving}.

\subsection{The Cayley/B\'ezout resultant}\label{sec:cayley}

\emph{Idea.} A divided difference compresses the common-root test into a
smaller matrix. Begin with two univariate polynomials $q_1,q_2$ of degree
at most $n$. Their \emph{B\'ezoutian} is
\begin{equation}\label{eq:bezoutian}
  F(s,t)=\frac{q_1(s)q_2(t)-q_2(s)q_1(t)}{s-t}
       =v_n(s)^TB(q_1,q_2)v_n(t),
  \qquad v_n(t)=(1,t,\ldots,t^{n-1})^T.
\end{equation}
The numerator vanishes on $s=t$, so the quotient is a polynomial of degree
at most $n-1$ in each variable. Interchanging $s,t$ changes both numerator
and denominator by a sign; hence $F(s,t)=F(t,s)$ and $B^T=B$. In particular,
this is complex symmetry, not Hermitian symmetry. If $\xi$ is a common
root, then $F(s,\xi)=0$ identically, and therefore $Bv_n(\xi)=0$.

\begin{proposition}[B\'ezout matrix and quotient multiplication]\label{rem:barnett}
  Suppose $q_1$ is monic of degree $n$, $\deg q_2\le n$, and $C_{q_1}$ is
  the subdiagonal companion matrix in \eqref{eq:companion}. Then
  \begin{equation}\label{eq:barnett}
    B(q_1,q_2)=q_2(C_{q_1})B(q_1,1)
             =B(q_1,1)q_2(C_{q_1}^{T}).
  \end{equation}
  Consequently
  \[
    \det B(q_1,q_2)=(-1)^{n(n-1)/2}
       \prod_{q_1(\alpha)=0}q_2(\alpha)
       =(-1)^{n(n-1)/2}\Res(q_1,q_2).
  \]
\end{proposition}

\begin{proof}
Put $H(s,t)=(q_1(s)-q_1(t))/(s-t)$. The identity
\[
  F(s,t)-q_2(s)H(s,t)
    =q_1(s)\frac{q_2(t)-q_2(s)}{s-t}
\]
shows that $F$ is the remainder of $q_2(s)H$ modulo $q_1(s)$.
Multiplication by $q_2(s)$ followed by this reduction acts on coefficient
columns as $q_2(C_{q_1})$, proving the first equality. Transpose it to obtain
the second. The entries of $B(q_1,1)$ on its anti-diagonal are $1$, and
those beyond it are zero, so its determinant is $(-1)^{n(n-1)/2}$.
Finally, the eigenvalues of $C_{q_1}$ are the roots of $q_1$, with
multiplicity, giving the determinant formula.
\end{proof}

This is Barnett's factorization; \cite{helmke1989bezoutians} gives a broader
account. Its order depends on the companion convention. If instead $q_1$
has degree $n$ and leading coefficient $a\ne0$, and $q_2$ has actual
degree $k\le n$, the formula becomes
$\det B=(-1)^{n(n-1)/2}a^{n-k}\Res(q_1,q_2)$.
The leading-coefficient factor therefore matters when the degrees differ.
Bilinearity in the two sets of input coefficients follows directly from
\eqref{eq:bezoutian}, rather than from powers of the companion matrix.

\begin{example}[The same roots in a smaller matrix]\label{ex:bezout-running}
  With $q_1(s)=s^2-1$, $q_2(s)=y^2-s$, direct division gives
  \[
    F(s,t)=-1+y^2(s+t)-st,\qquad
    B(y)=\begin{pmatrix}-1&y^2\\y^2&-1\end{pmatrix},\qquad
    \det B(y)=1-y^4.
  \]
  At $y=\eta$ with $\eta^4=1$, its kernel is spanned by $(1,\eta^2)^T$,
  which again gives $x=\eta^2$. The $2\times2$ quadratic matrix polynomial
  has a size-$4$ companion linearization, matching the four roots.
  The factorization is visible here: $C_{q_1}=B(q_1,1)=
  \left(\begin{smallmatrix}0&1\\1&0\end{smallmatrix}\right)$ and
  $B(y)=(y^2I-C_{q_1})B(q_1,1)$.
\end{example}

\emph{Several eliminated variables.} Write $q_j=p_j[x_d]$, viewing $x_d$
as a parameter, and introduce $s,t\in\C^{d-1}$. Define
\begin{equation}\label{eq:cayleyfun}
  F(s,t)=\frac{1}{\prod_{k=1}^{d-1}(s_k-t_k)}
  \det\begin{pmatrix}
    q_1(s_1,\ldots,s_{d-1})&\cdots&q_d(s_1,\ldots,s_{d-1})\\
    q_1(t_1,s_2,\ldots,s_{d-1})&\cdots&q_d(t_1,s_2,\ldots,s_{d-1})\\
    \vdots&&\vdots\\
    q_1(t_1,\ldots,t_{d-1})&\cdots&q_d(t_1,\ldots,t_{d-1})
  \end{pmatrix}.
\end{equation}
Successive rows agree when $s_k=t_k$, so every displayed factor divides
the determinant and $F$ is a polynomial. If each $q_j$ has degree at most
$n$ in each variable, $s_k$ occurs in the first $k$ rows and $t_k$ in the
last $d-k$ rows. Division therefore gives
\[
  \deg_{s_k}F\le kn-1,\qquad \deg_{t_k}F\le(d-k)n-1.
\]
Expand $F(s,t)=\sum_{\alpha,\beta}R_{\alpha\beta}s^\alpha t^\beta$ using
these full exponent boxes. Both index sets have size
$M=(d-1)!n^{d-1}$, producing a square matrix $R(x_d)$ of degree at most
$dn$ in the hidden variable. For $d>2$, this matrix need not be symmetric.

If $\xi$ is a common affine root of the specialized $q_j$, the last row in
\eqref{eq:cayleyfun} vanishes at $t=\xi$. Cancelling the nonzero polynomial
$\prod_k(s_k-\xi_k)$ proves $F(s,\xi)=0$. Similarly $F(\xi,t)=0$.
Thus $(\xi^\beta)_\beta$ and $(\xi^\alpha)_\alpha$ are right and transpose-left
null vectors. In a one-dimensional kernel they can be recovered up to scale.
A coordinate is read from an available pair of monomials differing by that
coordinate; coordinates absent from these vectors require back-substitution.

\emph{When is this an exact resultant?} For the full common box support,
the determinant is, up to a fixed nonzero normalization, the
multihomogeneous resultant of the $q_j$ on $(\mathbb P^1)^{d-1}$
\cite[Thm.~6.13]{dickenstein2003multihomogeneous}. This supplies the converse
to the null-vector argument, including boundary roots in that compactification.
For generic full-box $p_j$, the system has $d!n^d$ simple affine roots,
no boundary roots in $(\mathbb P^1)^d$, and distinct hidden coordinates.
Then $R(x_d)$ is regular
and its finite eigenvalues are exactly those coordinates.
Arbitrary sparse or total-degree inputs need a support-appropriate
construction: padding the full-box matrix can make it identically singular.

\begin{example}[A linear system and Cramer's rule]
  Take $p_1=x-z$, $p_2=y-z$, $p_3=x+y+z-1$ and hide $z$.
  Subtracting successive rows in \eqref{eq:cayleyfun} and dividing out the
  differences reduces its determinant to
  \[
    \det\begin{pmatrix}1&0&1\\0&1&1\\t_1-z&t_2-z&t_1+t_2+z-1\end{pmatrix}
    =3z-1.
  \]
  The support-adapted Cayley matrix is the scalar $3z-1$; back-substitution
  gives $x=y=z=1/3$. Its coefficients are precisely the determinants in
  Cramer's rule. Padding this constant-in-$s,t$ expression to the generic
  box size $2$ would introduce a zero row and column.
\end{example}

\begin{reduction}
  For generic full-box inputs, solve $R(x_d)v=0$ of size
  $(d-1)!n^{d-1}$ and degree $dn$, for example through a size-$d!n^d$
  linearization. Recover the other coordinates from null vectors or the
  original equations, checking boundary candidates for specialized inputs.
\end{reduction}

\emph{Numerical aspects and references.} Each entry is homogeneous of degree
$d$ in the input coefficients, and linear in the coefficients of each
polynomial separately. Scaling and the permitted perturbations therefore
matter when comparing this eigenproblem with the original root problem;
small matrix-coefficient errors are not automatically small input errors.
The classical divided-difference constructions are
\cite{cayley1857note,dixon1909eliminant}; multihomogeneous determinantal
formulae are treated in \cite{dickenstein2003multihomogeneous}.
See \cite{noferini2016numerical} for the numerical analysis, including the
general Cramer connection, and \cite{nakatsukasa2015computing} for a
bivariate implementation using subdivision and Chebyshev representations.
\subsection{The Macaulay resultant and the \texorpdfstring{$u$}{u}-resultant}\label{sec:macaulay}

\emph{Idea.} Adjoin a linear form whose vanishing tests a chosen projection of the roots.
If $F_i(x_0,x)$ is the homogenization of $p_i$ and
$h=\mu_0x_0+\sum_i\mu_ix_i$, the projective resultant
$\Res(F_1,\ldots,F_d,h)$ vanishes precisely when a common root of the $F_i$ lies on $h=0$.
With the $F_i$ fixed, this polynomial in $\mu$ is the \emph{$u$-resultant}.
When all projective roots are simple and lie in $x_0\ne0$, its product formula is
\[
  \Res(F_1,\ldots,F_d,h)=c\prod_{j=1}^N h(1,z_j),\qquad c\ne0.
\]
Thus a separating linear form turns the root set into a univariate spectrum
\cite{vanderwaerden1950modern,gelfand1994discriminants}.

\emph{A precise construction.} Put $e_i=\deg p_i$, and assume that the $F_i$ define a
reduced, zero-dimensional projective complete intersection with no point on $x_0=0$.
These are ordinary projective hypotheses; they need not hold for a generic system with
prescribed box support. Write
\[
  \widehat R=\C[x_0,\ldots,x_d],\qquad
  \widehat A=\widehat R/\langle F_1,\ldots,F_d\rangle,
  \qquad s=\sum_i(e_i-1),\quad \rho=s+1.
\]
Here $\widehat R_t$ and $\widehat A_t$ denote the homogeneous degree-$t$ pieces.
The complete-intersection
Hilbert series gives
\[
  \sum_{t\ge0}(\dim\widehat A_t)T^t
   =\frac{\prod_i(1+T+\cdots+T^{e_i-1})}{1-T},
  \qquad \dim\widehat A_t=N=\prod_i e_i\quad(t\ge s).
\]
Indeed, quotienting successively by a homogeneous nonzerodivisor of degree $e_i$
multiplies the Hilbert series by $1-T^{e_i}$; the displayed numerator has degree $s$
and value $N$ at $T=1$. Evaluation at $(1,z_j)$ identifies $\widehat A_s$ and
$\widehat A_\rho$ with
$\C^N$; this is the standard Hilbert-function description of a reduced projective
complete intersection \cite[Chs.~3 and~5]{cox2005using}.

Let $D_\rho$ have one column per degree-$\rho$ monomial and coefficient rows for all
$mF_i$ with $\deg m=\rho-e_i$. Its row space is precisely the degree-$\rho$ part of
$\langle F_1,\ldots,F_d\rangle$. Choose a matrix $Z$ whose $N$ columns form a basis of
$\ker D_\rho$, and choose homogeneous lifts $b_1,\ldots,b_N\in\widehat R_s$
of a basis of $\widehat A_s$.
Both choices can be made by rank-revealing linear algebra on polynomial coefficient
matrices. For a linear form $h$, form the $N\times N$ matrix
\begin{equation}\label{eq:macaulay-compression}
  H_h=\begin{pmatrix}
    \operatorname{coeff}_\rho(b_1h)\\[-1mm]\vdots\\
    \operatorname{coeff}_\rho(b_Nh)
  \end{pmatrix}Z.
\end{equation}
Coefficient vectors in this formula are rows, in the same monomial order as $D_\rho$.
This compresses the $h$-rows after imposing all polynomial relations.
Choose homogeneous linear forms $\alpha,\beta$ and put $h=\alpha-\lambda\beta$.

\begin{proposition}\label{prop:macaulay-compression}
  Under the stated assumptions, $\det H_h=c'\prod_jh(1,z_j)$ with $c'\ne0$ independent
  of $h$. If $\beta(1,z_j)\ne0$ for every $j$, then $H_\beta$ is invertible and the pencil
  $H_\alpha-\lambda H_\beta$ has eigenvalues
  $\lambda_j=\alpha(1,z_j)/\beta(1,z_j)$, counted with multiplicity.
\end{proposition}
\begin{proof}
Let $V_\rho$ have as its $j$th column the degree-$\rho$ monomials evaluated
at $(1,z_j)$, and put $E=(b_i(1,z_j))_{i,j}$. Evaluation is an isomorphism in these degrees,
so $E$ is invertible and $V_\rho$ has rank $N$. Since $D_\rho V_\rho=0$, we have
$Z=V_\rho C$ for some invertible $C$. Consequently
\[
  H_h=E\,\diag\bigl(h(1,z_1),\ldots,h(1,z_N)\bigr)C.
\]
Taking determinants proves the assertions.
\end{proof}

If the $\lambda_j$ are distinct, a right eigenvector $w_j$ satisfies
$Zw_j\propto V_\rho e_j$, so the root coordinates are the ratios of the entries indexed
by $x_ix_0^{\rho-1}$ and $x_0^\rho$. Taking $\beta=x_0$ is always valid here; a generic
$\alpha$ separates the roots. No explicit factorization of the $u$-resultant is needed.

\begin{example}[Macaulay compression for the running system]\label{ex:macaulay-running}
  For \eqref{eq:running}, $F_1=x^2-x_0^2$, $F_2=y^2-xx_0$, and $\rho=3$.
  There are ten degree-three monomials and six independent polynomial rows:
  $x_0F_1,xF_1,yF_1,x_0F_2,xF_2,yF_2$. Choose
  \[
    (b_1,b_2,b_3,b_4)=(x_0^2,x_0x,x_0y,xy),\qquad h=y-\lambda x_0.
  \]
  In degree three use the quotient basis $(x_0^3,x_0^2x,x_0^2y,x_0xy)$.
  Reduction by $x^2=x_0^2$ and $y^2=xx_0$ makes the four rows $b_ih$ equal to
  \[
    H_h=\begin{pmatrix}
      -\lambda&0&1&0\\
      0&-\lambda&0&1\\
      0&1&-\lambda&0\\
      1&0&0&-\lambda
    \end{pmatrix},\qquad \det H_h=\lambda^4-1.
  \]
  This is \eqref{eq:macaulay-compression} with $Z$ the coefficient-to-quotient map.
  At a root, $(1,x,y,xy)^T$ lies in its kernel when $\lambda=y$; its entries recover $x,y$.
  The constant matrix is the transpose of multiplication by $y$ in the affine basis
  $(1,x,y,xy)$, reflecting our use of coefficient rows.
\end{example}

\emph{Relation to the classical determinant formula.} Macaulay expresses the resultant
as a quotient of two specified determinants \cite{macaulay1902some,macaulay1916algebraic}.
An arbitrary square minor of a coefficient matrix need not equal the resultant or have
exactly its zero set. In particular, one cannot always retain all polynomial rows:
for three quadratics, $\rho=4$ and relations such as $F_2F_1-F_1F_2=0$ already make
those rows dependent. The compression above gives a direct, verifiable eigenproblem;
classical determinant quotients use their own prescribed row selection and extraneous-factor
removal. See J\'onsson--Vavasis \cite{jonsson2005accurate} for a numerical resultant construction.

\begin{reduction}
  An $N\times N$ generalized eigenvalue problem $H_\alpha w=\lambda H_\beta w$, obtained
  from homogeneous polynomial relations and a quotient basis. Under the hypotheses above
  it is regular and has exactly the projected roots as its eigenvalues.
\end{reduction}

\emph{Numerical aspects.} The uncompressed matrix has $\binom{\rho+d}{d}$ columns.
Computing its numerical rank, choosing a quotient basis, and scaling polynomial rows
are substantial parts of the algorithm. Algebraically equivalent compressions can have
different conditioning. The conditioning results for a particular Macaulay pencil
\cite{graf2026numerical} therefore need not transfer unchanged to every quotient compression.

\subsection{Multiparameter eigenvalue formulations}\label{sec:mep}

\emph{Idea.} Represent each polynomial as a determinant of an affine matrix pencil:
\[
  p_k(x)=\det W_k(x),\qquad
  W_k(x)=A_{k,0}+\sum_{i=1}^d x_iA_{k,i},\quad A_{k,i}\in\C^{m_k\times m_k}.
\]
Then $P(x)=0$ if and only if there are nonzero vectors satisfying the
\emph{multiparameter eigenvalue problem}
\begin{equation}\label{eq:mep}
  W_k(x)v_k=0,\qquad v_k\ne0,\qquad k=1,\ldots,d.
\end{equation}
The representation trades polynomial degree for matrix size. Explicit constructions
are available; a \emph{uniform} representation has entries affine-linear in the polynomial
coefficients as well as in the variables, allowing products of one coefficient and one
variable. For fixed $d\ge2$, the smallest size for a representation uniform over all
total-degree-at-most-$n$ polynomials is $\Theta(n^{d/2})$; in two variables explicit constructions have
size $2n+1$ or $2n-1$ \cite{boralevi2017uniform}. Choosing a representation is itself
an algorithmic step, separate from solving its eigenproblem.

\emph{Tensor Cramer's rule.} On $\C^{m_1}\otimes\cdots\otimes\C^{m_d}$ define
\[
  \Delta_0=\sum_{\sigma\in\mathfrak S_d}\operatorname{sgn}(\sigma)
      A_{1,\sigma(1)}\otimes\cdots\otimes A_{d,\sigma(d)},
\]
and define $\Delta_i$ by replacing column $i$ of the coefficient array
$(A_{k,l})_{k,l=1}^d$ with $(-A_{k,0})_k$ before this expansion.
For two parameters, writing $W_k=A_k+xB_k+yC_k$, this means
\begin{align*}
  \Delta_0&=B_1\otimes C_2-C_1\otimes B_2,\\
  \Delta_1&=C_1\otimes A_2-A_1\otimes C_2,\\
  \Delta_2&=A_1\otimes B_2-B_1\otimes A_2.
\end{align*}
If \eqref{eq:mep} holds and $v=v_1\otimes\cdots\otimes v_d$, then
\[
  \Delta_i v=x_i\Delta_0v,\qquad i=1,\ldots,d.
\]
To check this, substitute $-A_{k,0}v_k=\sum_lx_lA_{k,l}v_k$ into the determinant
expansion: every term with a repeated column cancels, leaving only the term $l=i$.
For scalar blocks this is ordinary Cramer's rule.

When $\Delta_0$ is invertible, the multiparameter problem is called \emph{nonsingular}.
Atkinson's theorem states that $\Gamma_i=\Delta_0^{-1}\Delta_i$ commute and that their
joint eigenvalues are exactly the solutions of \eqref{eq:mep}
\cite{atkinson1972multiparameter}. This converse is the substantive spectral theorem;
the preceding tensor calculation supplies the elementary forward implication.
One solves a separating linear combination $\sum_i a_i\Gamma_i$ and recovers all
coordinates on its common eigenvectors. If the chosen combination has repeated eigenvalues,
its invariant subspaces must be resolved using the remaining operators.

\begin{example}[A nonsingular two-parameter problem]\label{ex:mep-running}
  The running system has the representations
  \[
    W_1=\begin{pmatrix}x&1\\1&x\end{pmatrix},\qquad
    W_2=\begin{pmatrix}y&x\\1&y\end{pmatrix}.
  \]
  Their determinants are $x^2-1$ and $y^2-x$. In the tensor order
  $(e_1\otimes e_1,e_1\otimes e_2,e_2\otimes e_1,e_2\otimes e_2)$,
  \[
    \Delta_0=I_4,\qquad
    \Delta_1=\begin{pmatrix}0&0&-1&0\\0&0&0&-1\\-1&0&0&0\\0&-1&0&0\end{pmatrix},
    \quad
    \Delta_2=\begin{pmatrix}0&0&0&1\\-1&0&0&0\\0&1&0&0\\0&0&-1&0\end{pmatrix}.
  \]
  Direct multiplication gives $\Delta_2^2=\Delta_1$ and $\Delta_1^2=I_4$.
  Moreover $\det(\lambda I_4-\Delta_2)=\lambda^4-1$, so its four eigenvalues are
  $y=1,-1,i,-i$; the corresponding $x$ is $y^2$.
  At any root one can check the common eigenvector explicitly:
  \[
    v_1=(1,-x)^T,\quad v_2=(-y,1)^T,\quad
    v=v_1\otimes v_2=(-y,1,xy,-x)^T.
  \]
  Thus $\Delta_1v=xv$ and $\Delta_2v=yv$, and $x=-v_4/v_2$, $y=-v_1/v_2$.
  Solving only for $x$ leaves two-dimensional eigenspaces; $y$ separates all four roots.
\end{example}

\begin{reduction}
  Coupled generalized eigenproblems $\Delta_i v=x_i\Delta_0v$ on a tensor space of
  dimension $\prod_km_k$. In the nonsingular case, a separating linear combination
  gives an ordinary or generalized eigenproblem followed by recovery on common eigenvectors.
\end{reduction}

\emph{Numerical aspects.} Invertibility of $\Delta_0$ is a hypothesis to check, not a
consequence of having determinantal representations. A singular $\Delta_0$ does not by
itself make each pencil singular: a regular pencil can have eigenvalues at infinity.
Nor does excess tensor dimension alone distinguish infinite eigenvalues from a singular
pencil. For genuinely singular coupled pencils, a structure-aware staircase or related
reduction is needed before extracting the finite joint spectrum
\cite{plestenjak2016roots}. For the distinct issue of condition numbers of singular
matrix polynomials, see \cite{lotz2020wilkinson}. Representation size, scaling, and the tensor dimension all
affect practical accuracy and cost. Backward error and conditioning of the multiparameter
problem are studied in \cite{hochstenbach2003backward}; the conditioning of the
operator-determinant rootfinding reduction is studied in \cite{graf2026numerical}.
\section{Family II: scalar eliminants}\label{sec:familyII}

A scalar eliminant records the values of a coordinate or linear form on the root
set. It must be accompanied by a way to recover the other coordinates. The key
distinction is whether the chosen scalar value separates the roots.

\subsection{Gr\"obner elimination and the shape lemma}\label{sec:groebner}

For a coordinate $x_d$, the elimination ideal $I\cap\C[x_d]$ is generated by a
unique monic polynomial $f$. Proposition~\ref{prop:evaluation} shows directly that
\begin{equation}\label{eq:eliminant}
  f(T)=\prod_{\eta\in\{(z_j)_d:1\le j\le N\}}(T-\eta).
\end{equation}
Indeed, $h(x_d)$ belongs to $I$ exactly when $h$ vanishes at all these distinct values.
Thus $f$ is the \emph{minimal} polynomial of $M_{x_d}$; it equals the characteristic
polynomial only when $x_d$ separates the roots.

To compute $f$ from the equations, choose lexicographic order
$x_1\succ\cdots\succ x_d$. The elimination theorem says that the members involving
only $x_d$ of a Gr\"obner basis form a basis of $I\cap\C[x_d]$
\cite[Ch.~3]{cox2015ideals}. The short reason is that division of a univariate member
of $I$ must reduce its leading power of $x_d$; any Gr\"obner-basis leading monomial
dividing this power is itself a power of $x_d$, and lexicographic order then forces
the entire divisor to be univariate. Iterating gives a univariate reduction to zero.

\begin{proposition}[Shape lemma in the reduced case]\label{prop:shape}
If $x_d$ separates the roots, there are unique polynomials $h_i$ of degree less
than $N$ such that
\begin{equation}\label{eq:shape}
  I=\langle x_1-h_1(x_d),\ldots,x_{d-1}-h_{d-1}(x_d),f(x_d)\rangle,
  \qquad \deg f=N.
\end{equation}
The displayed generators form the reduced lexicographic Gr\"obner basis.
\end{proposition}

\begin{proof}
Interpolate $h_i((z_j)_d)=(z_j)_i$ at the $N$ distinct values. Each displayed generator
vanishes at all roots and hence belongs to $I$. The quotient by the displayed ideal
is isomorphic to $\C[T]/(f)$ and has dimension $N$, as does $R/I$, so the inclusion of
ideals is equality. The leading monomials $x_1,\ldots,x_{d-1},x_d^N$ are pairwise
coprime; the generators are monic and their other terms are reduced. This proves
the Gr\"obner-basis assertion by Buchberger's criterion.
\end{proof}

This elementary proof explains the usual shape lemma
\cite{gianni1989algebraic,cox2005using}: a separating coordinate makes the whole
quotient algebra generated by one element. A generic linear change of coordinates
provides such a coordinate, but does not guarantee that the projected roots are
well separated numerically.

\begin{example}[Elimination in the running system]\label{ex:shape}
With $x\succ y$, set $a=x-y^2=-q$. Reducing $p$ by $a$ gives
\[
  p-(x+y^2)a=y^4-1.
\]
Conversely, $p=(x+y^2)a+(y^4-1)$ and $q=-a$, so
\[
  I=\langle x-y^2,\ y^4-1\rangle.
\]
Solving $y^4-1=0$ and evaluating $x=y^2$ recovers all four roots. If we eliminate
$y$ instead, the eliminant is $x^2-1$, of degree two. Its two roots require the
additional fiber solves $y^2=x$; the characteristic polynomial of $M_x$ is
$(T^2-1)^2$. This is why an eliminant need not have degree $N$.
\end{example}

\begin{reduction}
Compute a lexicographic Gr\"obner basis, solve the univariate eliminant, and recover
the other coordinates by the shape formulas when the projection separates.
Without separation, solve the remaining fibers or change the projection. A
univariate companion solve adds its own representation and perturbation model.
\end{reduction}

\emph{Computation and literature.} Buchberger's algorithm and the F4/F5 algorithms
compute Gr\"obner bases over exact coefficient fields
\cite{buchberger2006algorithm,faugere1999f4,faugere2002f5}. With rational input one can
combine exact elimination with arbitrary-precision numerical rootfinding.
For approximate coefficients, deciding whether a leading coefficient is zero is a
rank decision; exact leading-term algorithms can change discontinuously across such
decisions. This motivates the numerical normal-form constructions of
Section~\ref{sec:normalform}. Coefficient growth and scalar root conditioning must
still be considered even when elimination itself is performed exactly
\cite{graf2026numerical}.

\subsection{Rational univariate representation}\label{sec:rur}

Choose a separating linear form $t=u^Tx$ and put $t_j=t(z_j)$. The reduced-case
rational univariate representation (RUR) consists of
\begin{equation}\label{eq:rur}
  f(T)=\prod_{j=1}^N(T-t_j),\qquad
  g_v(T)=\sum_{j=1}^N v(z_j)\prod_{k\ne j}(T-t_k),
\end{equation}
for $v=1,x_1,\ldots,x_d$. Here $g_1=f'$ and $\deg g_v<N$.

\begin{proposition}[Rational recovery and a coefficient construction]\label{prop:rur}
For every root $z_j$,
\[
  (z_j)_i=\frac{g_{x_i}(t_j)}{f'(t_j)}.
\]
The polynomials in \eqref{eq:rur} can be computed without knowing the roots by
\begin{equation}\label{eq:rur-trace}
  f(T)=\det(TI-M_t),\qquad
  g_v(T)=\tr\bigl(M_v\adj(TI-M_t)\bigr).
\end{equation}
\end{proposition}

\begin{proof}
At $T=t_j$, all terms in $g_v$ except the j-th vanish. The surviving product is
$f'(t_j)\ne0$, proving recovery. To prove \eqref{eq:rur-trace}, simultaneously
diagonalize $M_t$ and $M_v$ using \eqref{eq:stickelberger}; the adjugate of the
diagonal matrix has j-th entry $\prod_{k\ne j}(T-t_k)$, and trace gives the sum.
\end{proof}

The trace identity is a conceptual recipe. Implementations can obtain the needed
coefficients from traces $\tr(M_vM_t^k)$ and Newton identities, rather than forming
the full polynomial adjugate \cite{rouillier1999solving,alonso1996zeros}. In the
radical case, RUR and polynomial shape coordinates are equivalent: $f'$ is invertible
modulo $f$ and $h_i\equiv g_{x_i}/f'\pmod f$. The rational form is useful for exact
representation and arithmetic; it is not a general numerical conditioning guarantee.

\begin{example}[Rational recovery in the running system]\label{ex:rur}
Take $t=y$. We have $f(T)=T^4-1$, $f'(T)=4T^3$, and
\[
  g_x(T)\equiv T^2f'(T)\equiv4T\pmod f,\qquad
  g_y(T)\equiv Tf'(T)\equiv4\pmod f.
\]
At a root $\lambda$ of $f$ the recovered point is
\[
  x=\frac{4\lambda}{4\lambda^3}=\lambda^2,\qquad
  y=\frac4{4\lambda^3}=\lambda,
\]
using $\lambda^4=1$. A form need not be a coordinate: $t=x+y$ also separates
these roots, taking values $2,0,-1+i,-1-i$. Its polynomial is
\[
  f_t(T)=T(T-2)(T^2+2T+2)=T^4-2T^2-4T.
\]
Since $y=t-x$ and $x^2=1$, the equation $y^2=x$ gives
$x=(t^2+1)/(2t+1)$, with $y=t-x$; the denominator is nonzero at the four listed
values. This is an alternative rational parametrization of exactly the same roots.
\end{example}

\begin{reduction}
Find a separating form $t$; compute $f$ and the coordinate numerators; solve the scalar
degree-$N$ equation $f(T)=0$; evaluate the rational recovery formulas. If the chosen
form does not separate, these reduced-case formulas with denominator $f'$ cannot
be used unchanged.
\end{reduction}

\emph{Computation and literature.} Rouillier's RUR and Kronecker-style geometric
resolution make univariate descriptions effective
\cite{rouillier1999solving,giusti2001grobner,kronecker1882grundzuge}. An absolutely
continuous distribution for the coefficients of $t$ avoids the finitely many collision
hyperplanes with probability one. It need not avoid nearly colliding projections,
however, and both the scalar
solve and rational recovery can then be sensitive. Making $f$ monic fixes a scalar
normalization, not the coefficient norm, the scaling of $T$, or whether perturbations
may change its leading coefficient. Those choices remain part of a condition analysis.
\section{Family III: normal-form methods}\label{sec:familyIII}
\subsection{Constructing the multiplication operators}\label{sec:normalform}

Proposition~\ref{prop:multiplication} explains the eigenvalues. The computational
task is to obtain $[x_i b_k]_\B$ from the coefficients of $P$. A \emph{normal form}
is the linear reduction map taking a polynomial to its unique representative in
$\spanop\B$ modulo $I$. A Gr\"obner basis supplies one such map: repeatedly replace
leading monomials using the polynomial relations until only standard monomials
remain. Border bases allow more flexible choices of $\B$, with consistency conditions
ensuring that the resulting multiplication operators commute
\cite{mourrain1999criterion,elkadi2007introduction}.

\begin{example}[Multiplication matrices for the running system]\label{ex:toy}
The relations $x^2=1$, $y^2=x$ reduce every monomial to the span of
$\B=(1,x,y,xy)$. Since $\dim A=4$, these four classes form a basis. Columnwise
multiplication gives
\begin{equation}\label{eq:running-matrices}
M_x=\begin{pmatrix}
0&1&0&0\\1&0&0&0\\0&0&0&1\\0&0&1&0
\end{pmatrix},\qquad
M_y=\begin{pmatrix}
0&0&0&1\\0&0&1&0\\1&0&0&0\\0&1&0&0
\end{pmatrix}.
\end{equation}
For example, the last column of $M_y$ records $y(xy)=xy^2=x^2=1$.
Directly, $M_y^2=M_x$ and $M_y^4=I$. The vectors
$e_1,M_ye_1,M_y^2e_1,M_y^3e_1$ form a basis, so $M_y$ is cyclic with
minimal and characteristic polynomial $T^4-1$. Also, $M_x$ exchanges two pairs
of basis vectors. Thus
\[
\det(TI-M_y)=T^4-1,\qquad \det(TI-M_x)=(T^2-1)^2.
\]
Diagonalize $M_y$. If $M_y r=\lambda r$, then $M_xr=\lambda^2r$, so the
corresponding root is $(\lambda^2,\lambda)$. This gives all four roots with no
matching of independently ordered eigenvalue lists. Diagonalizing $M_x$ alone
would leave two two-dimensional eigenspaces that $M_y$ must further resolve.
\end{example}

Normal forms need only be known on a finite polynomial space containing $\B$ and
its products by the coordinates. This leads to the following practical formulation.

\begin{proposition}[A truncated normal form]\label{prop:tnf}
Let $U$ be a finite-dimensional polynomial space containing $\B$ and all $x_i\B$.
Suppose a linear map $L:U\to\C^N$ satisfies $\ker L=I\cap U$, and the classes
of $\B$ form a basis of $A$. Put
\[
  L_\B=(L(b_1)\ \cdots\ L(b_N)),\qquad
  L_{x_i\B}=(L(x_i b_1)\ \cdots\ L(x_i b_N)).
\]
Then $L_\B$ is invertible and
\begin{equation}\label{eq:tnf}
  M_{x_i}=L_\B^{-1}L_{x_i\B}.
\end{equation}
\end{proposition}

\begin{proof}
A linear combination of $\B$ annihilated by $L$ belongs to $I$, so its coefficients
vanish. This proves invertibility. Each column of the right side of
\eqref{eq:tnf} represents the unique combination of $\B$ whose difference from
$x_ib_k$ is annihilated by $L$, hence belongs to $I$.
\end{proof}

One constructs $L$ from linear relations between monomials: a sufficiently large
Macaulay matrix lists multiples of the input equations as rows. The quotient by
its row space, equivalently the cokernel of its transpose, gives the required
quotient map when the rank and degree conditions hold. The word
``sufficiently'' is essential; too small a truncation can miss relations.
QR with pivoting or an SVD can then select $\B$ for which $L_\B$ is well conditioned.
This is the approach of stabilized and truncated normal forms
\cite{telen2018stabilized,telen2018truncated,mourrain2021truncated}. It explains
both the flexibility of these methods and their connection to
Section~\ref{sec:macaulay}.

\begin{reduction}
Compute multiplication matrices in a chosen quotient basis. Choose a separating
$g=\sum_i\gamma_i x_i$ and solve the ordinary eigenproblem for
$M_g=\sum_i\gamma_i M_{x_i}$. Recover each coordinate from the action of the other
matrices on the same eigenspace. Equivalently, use evaluation eigenvectors of
$M_g^T$ when the basis contains appropriate coordinate monomials.
\end{reduction}

\emph{Numerical aspects and literature.} Multiplication matrices need not be normal,
although the running example happens to give permutation matrices. A Schur-based
implementation can work with clusters or invariant subspaces instead of individual
eigenvectors \cite{corless1997reordered}. The cost of a dense $N$-dimensional
eigensolve is $O(N^3)$; constructing a reliable quotient can itself be substantial.
Foundational numerical treatments are Auzinger--Stetter, M\"oller--Stetter, and
Stetter \cite{auzinger1988elimination,moller1995multivariate,stetter1996matrix,
stetter2004numerical}. For modern variants and the effects of clustering, see
\cite{telen2018truncated,parkinson2022analysis}. No basis-selection rule makes
intrinsically ill-conditioned roots disappear, and improving a quotient solve
does not by itself bound the conditioning of its eigenvectors.

\section{Conditioning: specify what is perturbed}\label{sec:conditioning}

The constructions above are exact algebraic statements. To compare their numerical
behavior, one must specify an input norm, an output norm, and which entries or
coefficients may change. This section gives the basic formulas rather than an
average-case classification.

\subsection{Equation values, coefficients, and representations}

At a simple root $z$, differentiation gives
\begin{equation}\label{eq:root-derivative}
  \delta z=-J(z)^{-1}\delta P(z).
\end{equation}
Thus $\|J^{-1}\|_2$ is the absolute condition for Euclidean perturbations of the
\emph{equation values}. Suppose instead that all equations use a fixed basis
$\phi=(\phi_1,\ldots,\phi_m)^T$, with $p_i=c_i^T\phi$, and the coefficient
array has its Euclidean norm.

\begin{proposition}[Coefficient condition of a root]\label{prop:coeff-condition}
With $S(z)=\sum_k|\phi_k(z)|^2$,
\[
  \kappa_{\mathrm{coeff}}(P,z)=\sqrt{S(z)}\,\|J(z)^{-1}\|_2.
\]
For a scalar output $g(z)$, replace $\|J^{-1}\|_2$ by $\|Dg(z)J^{-1}\|_2$.
\end{proposition}

\begin{proof}
The coefficient evaluation map $E_z:\delta c\mapsto\delta P(z)$ satisfies
$E_zE_z^*=S(z)I_d$. Therefore the squared operator norm of $J^{-1}E_z$ is
$S(z)\|J^{-1}\|_2^2$. Apply the chain rule for the scalar output.
\end{proof}

For example, with $p(T)=\sum_{k=0}^m a_kT^k$, the coefficient condition of a
simple root $\lambda$ is
\[
  \frac{\bigl(\sum_{k=0}^m|\lambda|^{2k}\bigr)^{1/2}}{|p'(\lambda)|},
\]
when all coefficients are perturbable. If a monic leading coefficient is held
fixed, the sum ends at $m-1$. The value-perturbation expression $1/|p'|$ omits
this evaluation norm. With a sup norm on a fixed unit disk, evaluation outside
the disk also needs care: a degree-$m$ perturbation bounded by $\epsilon$ on the
disk can reach $\epsilon|\lambda|^m$ at an exterior point.

For nonzero data and output, a normwise relative condition multiplies the
corresponding absolute condition by the input norm and divides by the output
norm. This does not change which perturbations are allowed. It also does not
remove dependence on a nonunitary change of basis.

\subsection{Eigenvectors, interpolation, and basis choice}

For a simple eigenvalue $\lambda$ of a matrix $M$, choose right and left
eigenvectors $r,w$, so $Mr=\lambda r$ and $w^*M=\lambda w^*$. Differentiating
and multiplying by $w^*$ gives
\[
  \delta\lambda=\frac{w^*(\delta M)r}{w^*r},\qquad
  \kappa_{\mathrm{eig}}=\frac{\|w\|_2\|r\|_2}{|w^*r|},
\]
for absolute spectral-norm matrix perturbations; a rank-one perturbation attains
the bound. For $M_g$ with $g$ separating, Proposition~\ref{prop:multiplication}
therefore gives
\begin{equation}\label{eq:overlap}
  \kappa_{\mathrm{eig}}(g(z_j),M_g)
  =\Lambda_j:=\|V^{-T}e_j\|_2\|Ve_j\|_2\ge1.
\end{equation}
The normalization is $v_j^TV^{-T}e_j=1$; in the Hermitian convention the left
eigenvector is $\overline{v_j}$.

The inverse Gram matrix makes the geometry explicit:
\[
  \Lambda_j^2=(V^*V)_{jj}\bigl((V^*V)^{-1}\bigr)_{jj}
  =\frac{\|v_j\|_2^2}
  {\operatorname{dist}(v_j,\spanop\{v_k:k\ne j\})^2}.
\]
To see the last equality, project $v_j$ orthogonally onto the complement of
the other columns. The vector dual to $v_j$ is this residual divided by its
squared norm. Thus conditioning measures an interpolation angle involving all
the evaluation vectors, not only the distance to the nearest root.

\begin{example}[A basis can create poor eigenvalue conditioning]\label{ex:basis}
For $p=x^2-1$, the basis $(1,x)$ gives
$M_x=\left(\begin{smallmatrix}0&1\\1&0\end{smallmatrix}\right)$, whose two
eigenvalues have condition one. In the basis $(1,\epsilon x)$, $\epsilon>0$,
\[
  M_x=\begin{pmatrix}0&\epsilon\\\epsilon^{-1}&0\end{pmatrix}.
\]
At the root 1 the evaluation vector is $(1,\epsilon)^T$ and the idempotent
$(1+x)/2$ has coefficient vector $(1/2,1/(2\epsilon))^T$. Hence
\[
  \Lambda=\frac{\epsilon+\epsilon^{-1}}2.
\]
The roots remain $\pm1$ and $|p'(\pm1)|=2$, but the unstructured eigenvalue
condition grows without bound as $\epsilon\to0$. This is a representation effect.
\end{example}

\subsection{Matrix polynomials and scalar eliminants}

For a regular matrix polynomial $R(\lambda)=\sum_{k=0}^mR_k\lambda^k$ and a
simple finite eigenvalue, let $R(\lambda)v=0$ and $w^*R(\lambda)=0$.
The same differentiation gives
\[
  \delta\lambda=-\frac{w^*\delta R(\lambda)v}{w^*R'(\lambda)v}.
\]
Under the coefficient norm $(\sum_k\|\delta R_k\|_F^2)^{1/2}$, its condition is
\begin{equation}\label{eq:pep-condition}
  \frac{\bigl(\sum_{k=0}^m|\lambda|^{2k}\bigr)^{1/2}
  \|w\|_2\|v\|_2}{|w^*R'(\lambda)v|}.
\end{equation}
Cauchy--Schwarz proves the upper bound, and coefficientwise aligned rank-one
perturbations attain it. Other coefficient weights or function norms give different
evaluation factors \cite{tisseur2000backward}. For the pencil $(M,I)$, allowing
independent perturbations of both matrices with their maximum spectral norm gives
$\Lambda(1+|\lambda|)$; holding $I$ fixed gives \eqref{eq:overlap}.

For a monic scalar eliminant with distinct roots $t_j$,
\[
  f'(t_j)=\prod_{k\ne j}(t_j-t_k).
\]
Taking logarithms expresses its value-perturbation sensitivity through pairwise
projected distances. This elementary identity does not by itself compare coefficient
conditioning across methods: the evaluation factors, eliminant coefficient norm,
and coordinate-recovery step still enter. A projection can bring distinct,
well-separated multivariate roots close together.

Finally, if $R(c,\lambda)$ is built exactly from the input coefficients, then a
coefficient perturbation induces $\delta R=D_cR[\delta c]$, a structured subset
of all matrix-polynomial perturbations. The chain rule must recover
\eqref{eq:root-derivative} for the exact composite solve and recovery map.
A backward-stable QR or QZ computation controls perturbations of the matrix or
pencil it receives; translating them back to the original polynomial coefficients
requires additional analysis of construction, linearization, and recovery.

Noferini--Townsend and Graf--Townsend exhibit systems for which intermediate
subproblems are much more sensitive than the original roots under their specified
perturbation models \cite{noferini2016numerical,graf2026numerical}. These are
worst-case results, not statements that every input is unstable. Nor do the exact
identities in this survey establish favorable average-case behavior. They identify
the quantities and choices that a probabilistic or smoothed analysis must control.

\section{Synthesis and comparison}\label{sec:synthesis}

\subsection{What the representations have in common}

The quotient algebra gives a precise comparison. For a separating $t$, the
eliminant is $\det(TI-M_t)$; without separation, its squarefree minimal polynomial
has smaller degree. Under the stated regularity and no-extraneous-root assumptions,
a hidden-variable resultant determinant represents the same projected root values.
Macaulay quotient compression and truncated normal forms compute different
coordinates for the same finite-dimensional multiplication maps.

One should therefore distinguish equality of the represented zero set, equality of
scalar determinants including multiplicities and factors, and similarity or equivalence
of matrix representations. The first does not imply the latter two, and none alone
implies equality of numerical condition numbers.

\subsection{An optional duality viewpoint: residues}\label{sec:residues}

There is a useful explanation of the Jacobian factors that appear with B\'ezoutians.
For a reduced complete intersection, define
\[
  \tau(f)=\tr M_f=\sum_j f(z_j),\qquad
  \mathcal R(f)=\sum_j\frac{f(z_j)}{\det J(z_j)}.
\]
The latter is the global residue functional in the simple-root case
\cite{cattani2005residues}. The pairings $\tau(fg)$ and $\mathcal R(fg)$ are
nondegenerate, but they are different: if $\mathcal J=[\det DP]\in A$, then
$\tau(f)=\mathcal R(\mathcal J f)$. Nondegeneracy follows immediately in the
idempotent basis, where both pairings are diagonal with nonzero diagonal entries.

\begin{proposition}[Univariate B\'ezoutian and residue duality]\label{prop:residue-bezout}
Let $p$ be monic and squarefree of degree $m$. Then $B(p,1)=H^{-1}$, where $H$ is the
residue Gram matrix in the monomial basis of $\C[x]/(p)$:
\[
  H_{ab}=\sum_{p(\lambda)=0}\frac{\lambda^{a+b}}{p'(\lambda)},
  \qquad 0\le a,b<m,
\]
\end{proposition}

\begin{proof}
For each root $\lambda_j$, the B\'ezoutian evaluated at its second argument is
$p(s)/(s-\lambda_j)=p'(\lambda_j)\ell_j(s)$. Thus
$B(p,1)V=V^{-T}\diag(p'(\lambda_j))$, whereas
$H=V\diag(1/p'(\lambda_j))V^T$. Multiplication gives $B(p,1)H=I$.
\end{proof}

For $p=x^2-1$, this residue Gram matrix and its inverse are both
$\left(\begin{smallmatrix}0&1\\1&0\end{smallmatrix}\right)$; the ordinary
trace Gram matrix is $2I$. Calling B\'ezoutian duality simply trace duality would
miss this Jacobian weight.

More generally, if $P(x)=Q(x)(x-z_j)$, then
\[
  [\det Q]=\det J(z_j)[\ell_j]\quad\text{in }A.
\]
At another root $z_k$, $Q(z_k)(z_k-z_j)=0$ makes its determinant vanish; at
$z_j$, differentiation gives $Q(z_j)=J(z_j)$. Evaluation proves the identity.
It explains why determinant numerators in normal-form condition formulas can
equally be written as norms of Lagrange idempotents \cite{graf2026numerical}.

\subsection{A comparison to use when choosing a construction}

\begin{table}[htbp]
\centering\small
\begin{tabularx}{\textwidth}{@{}>{\raggedright\arraybackslash}p{0.19\textwidth}XX@{}}
\toprule
Method & Intermediate problem and size & Conditions and recovery \\
\midrule
Sylvester & For two generic box-degree-$n$ equations: a $2n\times2n$ PEP of grade $n$.
 & Distinct projected roots and no common fiber root at infinity; recover the
 eliminated coordinate from a null vector or a fiber solve. \\[0.5em]
Cayley--B\'ezout & For $d$ generic box-degree-$n$ equations: size
 $(d-1)!n^{d-1}$, grade dn.
 & Regularity and the stated multihomogeneous model; recover coordinates from
 tensor evaluation vectors. \\[0.5em]
Macaulay / $u$-resultant & Dense degree-$\rho$ relation matrix, followed here by
 an $N\times N$ compressed pencil.
 & Complete intersection with no ordinary projective root at infinity; valid
 quotient ranks, a nonvanishing denominator, and a separating ratio
 $\alpha/\beta$ for individual eigenvector recovery. \\[0.5em]
Multiparameter & Operator-determinant pencils on a tensor space of dimension
 $\prod_k m_k$.
 & The nonsingular theory assumes $\Delta_0$ invertible. Singular formulations
 need an appropriate reduction; coordinates belong to common eigenspaces. \\[0.5em]
Gr\"obner elimination & A scalar eliminant of degree at most $N$, plus extension
 equations.
 & A separating coordinate gives degree $N$ and polynomial shape recovery;
 otherwise solve the remaining fibers. \\[0.5em]
RUR & A scalar degree-$N$ polynomial and $d$ rational coordinate formulas.
 & A separating linear form and nonvanishing recovery denominator; construction
 is symbolic or based on quotient traces. \\[0.5em]
Normal form & An $N\times N$ ordinary multiplication eigenproblem.
 & A valid quotient basis or truncation and a separating multiplier; recover
 coordinates on shared eigenspaces. \\
\bottomrule
\end{tabularx}
\caption{The seven constructions. A PEP is a polynomial eigenvalue problem.
Sizes in the first two rows refer to full box supports; the Macaulay row instead
uses total degrees and $\rho=\sum_i(e_i-1)+1$. A regular size-$m$ PEP of grade $r$
has $mr$ eigenvalues counting infinity and algebraic multiplicity; a companion
linearization has size $mr$. These counts are not guarantees of distinct affine roots.}
\label{tab:methods}
\end{table}

For exact input and symbolic output, elimination and RUR provide explicit root
descriptions. For numerical work, retaining matrix structure can avoid forming
large scalar eliminant coefficients, but introduces basis, linearization, and
eigenvector issues. Normal forms make basis selection explicit; resultants make
the relation to coefficient elimination explicit; multiparameter formulations make
affine matrix pencils explicit. The right choice depends on the polynomial support,
available arithmetic, desired output, and the conditioning of the construction.

The running example gives a concrete cross-check: Sylvester yields $y^4-1$,
B\'ezout yields $1-y^4$, Gr\"obner elimination gives $x=y^2$ and $y^4=1$, RUR
gives the same coordinates rationally, and multiplication by $y$ has characteristic
polynomial $T^4-1$. The Macaulay and multiparameter examples recover the same four
points through pencils. Their agreement is an algebraic fact; their response to
floating-point perturbations depends on which of these representations is used.

\clearpage
\bibliographystyle{plain}
\bibliography{survey}
\end{document}
